\documentclass[10pt,a4paper]{amsart}
\usepackage[margin=19mm]{geometry}
\usepackage{amsmath,amssymb,mathtools}
\usepackage[scr=rsfs,cal=euler]{mathalfa}
\usepackage{xfrac}
\usepackage[unicode,hidelinks,bookmarks=false]{hyperref}
\newcommand{\ie}{i.e.\ }
\newcommand{\cf}{cf.\ }
\newcommand{\resp}{respectively\ }
\newcommand{\Subsection}{\subsection}
\providecommand{\nobreakdash}{\nobreak\hbox{-}}
\setlength{\emergencystretch}{2em}
\multlinegap0pt
\usepackage{bm}
\usepackage{bbm}
\usepackage{dsfont}
\usepackage{xspace}
\usepackage{cancel}
\usepackage{mathtools}
\usepackage{amsthm}
\makeatletter
\@ifundefined{proposition}{\newtheorem{proposition}{Proposition}}{}
\@ifundefined{theorem}{\newtheorem{theorem}{Theorem}}{}
\makeatother
\makeatletter
\expandafter\ifx\csname Abstract\endcsname\relax
\newenvironment{Abstract}
{
\begin{center}
\textbf{Abstract}\\
\vspace{0.25cm}
\begin{minipage}{14.5cm}}
{\footnotesize
\end{minipage}
\end{center}}
\fi
\newcommand{\dd}{\operatorname{d}\,}
\expandafter\ifx\csname keywords\endcsname\relax
\newenvironment{keywords}
{
\begin{center}
\textbf{Keywords}\\
\vspace{0.17cm}
\begin{minipage}{14.5cm}}
{\footnotesize
\end{minipage}
\end{center}}
\fi
\makeatother
\title[A description of classical field equations using extensions ]{A description of classical field equations using extensions of graded Poisson brackets}
\author{Manuel de Le\'on, Rub\'en Izquierdo-L\'opez}
\date{Source: arXiv:2507.04743v1 (CC BY 4.0)}
\begin{document}
\maketitle

\noindent\emph{Source and licence.} A description of classical field equations using extensions of graded Poisson brackets. Version arXiv:2507.04743v1 (CC BY 4.0), distributed under the Creative Commons Attribution 4.0 International licence, as recorded in the arXiv licence metadata for this exact version and shown on the abstract page https://arxiv.org/abs/2507.04743v1. Full rights notice: licenses/arxiv-2507.04743-CC-BY-4.0.txt.\par\smallskip

\noindent\textit{Editorial scope.} Counted unit: the Proposition whose source label is \texttt{prop:sharp\_of\_bracket} in the article (source0.tex lines 868--874, source label \texttt{prop:sharp\_of\_bracket}), delivered together with the article's own complete original proof of it (source0.tex lines 888--919, one \texttt{proof} environment of 32 lines). No mathematics is written, repaired or re-derived: statement and proof are the article's own text line for line. The proof is detached from the statement in the source: the article states the result at lines 868--874 and prints its proof at lines 888--919, and the proof environment's own header names the statement, so the association is the article's own. Required context retained uncounted: thm:extension\_of\_sharp\_1 (lines 734--745). Every definition, display and label of those lines is retained unchanged. Omitted from this package and not redistributed: 1--733, 746--867, 875--887, 920--1910. Nothing inside a retained excerpt is omitted or altered: every retained body is delivered exactly as the article's lines are written, so no omission receipt of any kind accompanies this package. Citations: the retained excerpts carry no citation command at all, so no bibliography entry of the article is transcribed and no reference is redistributed. Reference adaptation: none was needed and none was made; every reference inside the delivered unit resolves inside it, so no unresolved reference or citation is produced. Printed slips kept literal: no printed word, symbol, hypothesis or number of the counted statement or of its proof was altered. Layout and class: the article's own class is \texttt{article} with packages including graphicx, geometry, amsfonts, parskip, amsthm, dsfont, fontenc, amsmath, mathtools, biblatex, enumitem, hyperref, tikz, xcolor; the wrapper is the lane's pinned \texttt{amsart} editorial wrapper at 10pt on A4, which restates the theorem-like environments the retained text needs and adds the article's own macro definitions that the retained text needs (\texttt{\textbackslash dd}), with extra wrapper packages (bm, bbm, dsfont, xspace, cancel, mathtools). The delivered numbering is the unit's own. Assets: no retained span contains a figure environment, a graphics-inclusion command, a PDF-page inclusion, a table or a TikZ picture; no image file is copied into this package, the package declares no asset dependency and no image file is redistributed with it. Licence: version arXiv:2507.04743v1 of this article is distributed under the Creative Commons Attribution 4.0 International licence, as recorded in the arXiv licence metadata for this exact version and shown on the abstract page https://arxiv.org/abs/2507.04743v1. A full rights notice is delivered at licenses/arxiv-2507.04743-CC-BY-4.0.txt. Characters: every retained excerpt is pure ASCII, so the delivered engine, XeLaTeX with its default Latin Modern text font, reports no missing character.
\begin{theorem}
\label{thm:extension_of_sharp_1}
Let $(M, S^a, \sharp_a)$ be a regular graded Dirac manifold of order $n$. Then, there exists an unique extension of $\sharp_1$, namely,
\[
(S^1)^{\wedge a} \xrightarrow{\widetilde \sharp_1} \bigwedge^{a - 1}M \otimes_M \bigvee_n M/K_n\,,
\]
for all possible values of $a \geq 1$, such that for any $\alpha \in S^n$ and $\beta \in (S^1)^{\wedge a}$ we have
\[
\iota_{\sharp_n(\alpha)} \beta  = (-1)^{(n + 1 - a)} \iota_{\widetilde \sharp_1(\beta)} \alpha\,.
\]

\end{theorem}

\begin{proposition}
\label{prop:sharp_of_bracket}
Let $\Theta \in \Omega^{a}_H(M)[1],$ and $\alpha \in \Omega^{n-1}_H(M).$ Then,
\[
\widetilde \sharp_1(\dd \{\alpha, \Theta\}) = -\pounds_{\sharp_n (\dd \alpha)} \left(\widetilde{\sharp}_1(\dd \Theta)  \right)\,. 
\]
\end{proposition}

\begin{proof}[Proof of Proposition \ref{prop:sharp_of_bracket}] In order to prove the equality, let us check that contraction against an arbitrary form $\beta$ taking values in $S^n$ by both sides of the desired equality yields the same result. Indeed, using the defining property of $\widetilde \sharp_1$ (see Theorem \ref{thm:extension_of_sharp_1}) and that $\dd \Theta \in S^{n+1}$,  we get
\begin{align*}
    \iota_{\widetilde \sharp_1(\dd \{\alpha, \Theta\})} \beta &= (-1)^{n-a} \iota_{\sharp_n(\beta)} \dd \{\alpha, \Theta\} = - \iota_{\sharp_n(\beta)} \dd \iota_{\widetilde \sharp_1(\dd \Theta) }\dd \alpha\\
    &= - \pounds_{\sharp_n(\beta)} \iota_{\widetilde \sharp_1(\dd \Theta)} \dd \alpha + \dd \iota_{\sharp_n(\beta)} \iota_{\widetilde \sharp_1(\dd \Theta)} \dd \alpha\\
    &= - (-1)^{(n-a)}\pounds_{\sharp_n(\beta)} \iota_{ \sharp_n(\dd \alpha)} \dd \Theta + \dd \iota_{\sharp_n(\beta)} \iota_{\widetilde \sharp_1(\dd \Theta)} \dd \alpha\,.
\end{align*}
Now, rewriting the leftmost term employing the formula $\pounds_X \iota_Y \omega  = \iota_{[X, Y]}\omega + \iota_Y \pounds_X \omega$, we get
\begin{align}
    \iota_{\widetilde \sharp_1(\dd \{\alpha, \Theta\})} \beta &= - (-1)^{n-a} \left(\iota_{[\sharp_n(\beta), \sharp_n(\dd \alpha)]} \dd \Theta + \iota_{\sharp_n(\dd \alpha)}  \pounds_{\sharp_n(\beta)} \dd \Theta \right)+\dd \iota_{\sharp_n(\beta)} \iota_{\widetilde \sharp_1(\dd \Theta)} \dd \alpha\,.
    \label{eq:contraction_unsing_Lie_bracket}
\end{align}
Recall that $(S^a, \sharp_a)$ is a graded Dirac structure, so that we have
$[\sharp_n(\beta), \sharp_n(\dd \alpha)] = \sharp_n(\eta)$, where 
\begin{equation}
    \eta = \pounds_{\sharp_n(\beta)} \dd \alpha - \pounds_{\sharp_n(\dd \alpha)} \beta + \frac{1}{2} \dd \left( \iota_{\sharp_n(\dd \alpha)} \beta - \iota_{\sharp_n(\beta)} \dd \alpha\right) = - \pounds_{\sharp_n(\dd \alpha)} \beta\,.
    \label{eq:description_of_eta}
\end{equation}
Using Eq. \eqref{eq:contraction_unsing_Lie_bracket} together with Eq. \eqref{eq:description_of_eta} we obtain
\begin{align*}
    \iota_{\widetilde \sharp_1(\dd \{\alpha, \Theta\})} \beta &=- (-1)^{n-a} \left(\iota_{[\sharp_n(\beta), \sharp_n(\dd \alpha)]} \dd \Theta + \iota_{\sharp_n(\dd \alpha)}  \pounds_{\sharp_n(\beta)} \dd \Theta \right)+\dd \iota_{\sharp_n(\beta)} \iota_{\widetilde \sharp_1(\dd \Theta)} \dd \alpha\\
    &= - (-1)^{n-a} \left( \iota_{\sharp_n(\eta)}  \dd \Theta + \iota_{\sharp_n(\dd \alpha)}  \pounds_{\sharp_n(\beta)} \dd \Theta\right)+\dd \iota_{\sharp_n(\beta)} \iota_{\widetilde \sharp_1(\dd \Theta)} \dd \alpha\\
    &= - (-1)^{n-a} \left( (-1)^{n-a} \iota_{\widetilde \sharp_1(\dd \Theta)} \eta + \iota_{\sharp_n(\dd \alpha)}  \pounds_{\sharp_n(\beta)} \dd \Theta\right) +\dd \iota_{\sharp_n(\beta)} \iota_{\widetilde \sharp_1(\dd \Theta)} \dd \alpha\\
    &=  \iota_{\widetilde \sharp_1(\dd \Theta)} \pounds_{\sharp_n(\dd \alpha)} \beta - (-1)^{n-a} \iota_{\sharp_n(\dd \alpha)} \pounds_{\sharp_n(\beta)} \dd \Theta + \dd \iota_{\sharp_n(\beta)} \iota_{\widetilde \sharp_1(\dd \Theta)} \dd \alpha \,.
\end{align*}
Now it only remains to rewrite $\iota_{\widetilde \sharp_1(\dd \alpha)} \pounds_{\sharp_n(\dd \alpha)} \beta =  -\iota_{\pounds_{\sharp_n(\dd \alpha)}\widetilde \sharp_1(\dd \Theta)} \beta +\pounds_{\sharp_n(\dd \alpha)}\iota_{\widetilde \sharp_1 (\dd \Theta)} \beta$ to get
\begin{align*}
    \iota_{\widetilde \sharp_1(\dd \{\alpha, \Theta\})} \beta =&- \iota_{\pounds_{\sharp_n(\dd \alpha)}\widetilde \sharp_1(\dd \Theta)} \beta + \pounds_{\sharp_n(\dd \alpha)}\iota_{\widetilde \sharp_1 (\dd \Theta)} \beta\\
    & - (-1)^{n-a} \iota_{\sharp_n(\dd \alpha)} \pounds_{\sharp_n(\beta)} \dd \Theta + \dd \iota_{\sharp_n(\beta)} \iota_{\widetilde \sharp_1(\dd \Theta)} \dd \alpha\\
    &= - \iota_{\pounds_{\sharp_n(\dd \alpha)}\widetilde \sharp_1(\dd \Theta)} \beta\,,
\end{align*}
finishing the proof.
\end{proof}

\end{document}
